\documentclass[11pt,letterpaper]{article}

\usepackage[utf8]{inputenc}
\usepackage[T1]{fontenc}
\usepackage[left=0.7in,right=0.7in,top=0.6in,bottom=0.6in]{geometry}
\usepackage{enumitem}
\usepackage{titlesec}
\usepackage{hyperref}
\usepackage{xcolor}
\usepackage{graphicx}
\usepackage{helvet}
\renewcommand{\familydefault}{\sfdefault}

\usepackage{microtype}
\usepackage[none]{hyphenat}
\sloppy
\frenchspacing

% Match font and spacing style from CV (no indent, standard Computer Modern, fixed gap)
\setlength{\parindent}{0pt}
\setlength{\parskip}{1ex}

\hypersetup{
    colorlinks=true,
    linkcolor=black,
    filecolor=black,      
    urlcolor=blue,
    citecolor=black,
}

\pagestyle{empty}

\begin{document}

\begin{center}
{\LARGE \textbf{Letter of Recommendation for Xuan Thanh Trinh}}\\[1.5ex]
{\small School of Electrical and Electronic Engineering, HUST, Hanoi, Vietnam $\vert$ son.tranthanh@hust.edu.vn}
\end{center}
\vspace{-3ex}
\noindent\rule{\textwidth}{0.4pt}
\vspace{-2ex}

\noindent Viet Nam,July 22, 2026

\vspace{0.2ex}

\noindent
\textbf{To:} Ph.D. Selection Committee \\
Doctoral Program in Electrical, Electronics and Computer Engineering \\
College of Engineering and Computer Science (CECS) \\
University of Michigan--Dearborn, MI 48128, USA

\vspace{0.5ex}

\noindent Dear Admissions Committee,

I am Tran Thanh Son, a Lecturer in the School of Electrical and Electronic Engineering at Hanoi University of Science and Technology (HUST). As a senior researcher in power system optimization and microgrid architectures (Ph.D., Shibaura Institute of Technology), I served as Xuan Thanh Trinh's direct research mentor on a complex project regarding peer-to-peer energy trading frameworks. Based on his exceptional performance, I firmly advocate for his admission to the Ph.D. program at the University of Michigan-Dearborn

Thanh demonstrates a rare capacity for mathematical maturity and autonomous research execution. Rather than passively applying formulas, he rigorously evaluates their fundamental analytical properties before deployment. Notably, he single-handedly established the global convergence proof for an ADMM framework applied to a Second-Order Cone Programming (SOCP) relaxation, which effectively convexifies complex, non-linear grid constraints. Operating with high autonomy, he consistently navigates dense academic literature to master theoretical frameworks outside the undergraduate curriculum, requiring minimal direct supervision. This dual capability profound mathematical rigor coupled with exceptional self-direction—proves his readiness for independent Ph.D. research.

Furthermore, Thanh exhibits exceptional network perspective and comprehensive framework validation. Specifically, he rejected the standard simplification of microgrids as passive loads and directly modeled the complex, bidirectional energy exchanges between autonomous microgrids and the distribution grid under high renewable penetrations to capture true network interactions. To verify the robustness of his framework, he demonstrated critical methodological rigor when he challenged his system formulations with severe stress testing across complex radial and meshed topologies, including the IEEE 123-bus benchmark, completely avoiding simplified test cases that only ensure easy convergence. By successfully aligning abstract formulations with practical grid constraints, he generated the core numerical validation for our manuscript, currently under review at Sustainable Energy, Grids and Networks (SEGAN).

Ultimately, Thanh's methodological resilience sets him apart during complex simulations. When debugging large-scale grid networks, he constructively treats algorithmic non-convergence as diagnostic data rather than failures. Even under strict deadlines, he systematically traces numerical instabilities to their theoretical root causes rather than relying on unverified heuristics. This analytical endurance proves he is exceptionally well-equipped for the exacting demands of a Ph.D. program.

In summary, Thanh uniquely combines mathematical rigor, network perspective, and methodological resilience. As his mentor and co-author, I give him my highest recommendation for admission to your Ph.D. program without reservation. I am confident he will be an immediate and highly productive contributor to your research group. Please contact me if you require further information.

\noindent Sincerely,

\vspace{1.9cm}

\noindent \textbf{Dr. Tran Thanh Son}\\
\noindent Lecturer, School of Electrical and Electronic Engineering\\
\noindent Hanoi University of Science and Technology (HUST), Vietnam
% \noindent Email: son.tranthanh@hust.edu.vn
\end{document}
